\documentclass[10pt,twocolumn,letterpaper]{article}

\usepackage{graphicx}
\usepackage{amsmath}
\usepackage{amssymb}
\usepackage{booktabs}
\usepackage[table,xcdraw]{xcolor}
\usepackage{listings}
\usepackage{pgfplots}
\usepackage{pgfplotstable}
\usepackage{tikz}
\usetikzlibrary{positioning}
\usepackage{subcaption}
\usepackage{url}
\usepackage[numbers]{natbib}
\usepackage{hyperref}
\usepackage[capitalize]{cleveref}

\pgfplotsset{compat=1.18}

\lstset{
  basicstyle=\ttfamily\footnotesize,
  breaklines=true,
  frame=single,
  keywordstyle=\bfseries,
  commentstyle=\itshape,
  showstringspaces=false
}

\begin{document}

\title{Code Complexity Metrics for AI-Assisted Development:\\
A Survey and Comparative Tool Analysis}

\author{Brandon Arrendondo\\
{\tt\small brandon.arrendondo@bissell.com}}

\maketitle

\begin{abstract}
Code complexity metrics have guided software quality practice for nearly
five decades, from McCabe's cyclomatic complexity~\cite{mccabe} in 1976
through Cognitive Complexity~\cite{cognitive-complexity} in 2018.  These
metrics were designed to model \emph{human} cognitive load.  The rise of
AI-assisted development introduces a different cost model: the effort and
token budget an LLM requires to safely modify a function.  This survey
reviews 10 established complexity metrics, evaluates 8 open-source
analysis tools across 5 dimensions (metric coverage, language support,
CI integration, output formats, and speed), and characterizes the gap
between existing metrics and the requirements of AI-cost modeling.  No
surveyed tool computes metrics oriented toward LLM modification cost, and
only two support the multi-language, CI-native deployment model needed for
modern polyglot codebases.  We describe the design requirements for an
AI-cost-aware complexity tool and position knots~\cite{knots} as the first
open-source tool to address this gap.
\end{abstract}

%% ============================================================
\section{Introduction}
\label{sec:intro}

Software complexity metrics serve two related purposes: they identify
functions that are difficult to understand, and they provide actionable
signals for refactoring.  Both purposes assume a model of the reader.
McCabe's cyclomatic complexity models the \emph{structural} cost of
testing all paths through a function.  Cognitive Complexity models the
\emph{reading} difficulty a human programmer experiences.  Halstead's
software science models \emph{vocabulary and volume} as proxies for
mental effort.

None of these metrics was designed for the reader a modern development
team uses most: an LLM assistant such as Claude, GPT-4, or Gemini.
LLM-assisted development has a fundamentally different cost model.  The
cost of asking an LLM to modify a function is not measured in programmer
minutes; it is measured in:

\begin{enumerate}
    \item \textbf{Reasoning effort} --- the number of branch paths, nesting
        levels, and state variables the model must track to generate a
        correct edit.
    \item \textbf{Context tokens} --- the tokens consumed loading the
        function itself and all external symbols it references before
        reasoning can begin.
\end{enumerate}

Traditional complexity metrics partially correlate with these costs but
were not calibrated against them.  A 300-line function with simple linear
structure scores high on SLOC but low on McCabe; it is cheap to reason
about (few branches) but expensive to tokenize (many lines).  A 20-line
function with 15 external calls scores low on all traditional metrics
but has high context pressure: the model must resolve 15 dependency
chains before acting.

This survey makes three contributions:

\begin{enumerate}
    \item A structured review of 10 established complexity metrics,
        organized by what aspect of complexity each models.
    \item A comparative evaluation of 8 open-source analysis tools on
        5 dimensions relevant to CI-native deployment.
    \item A characterization of the gap between existing tools and
        AI-cost-aware complexity analysis, motivating the AIRD and AICP
        metrics described in the companion paper~\cite{knots-aird}.
\end{enumerate}

The rest of this paper is organized as follows.
\Cref{sec:metrics} surveys established complexity metrics.
\Cref{sec:tools} evaluates open-source tools.
\Cref{sec:gap} characterizes the gap.
\Cref{sec:related} covers related work.
\Cref{sec:conclusion} concludes.

%% ============================================================
\section{Complexity Metrics}
\label{sec:metrics}

We organize complexity metrics into four categories: structural,
size-based, vocabulary-based, and AI-cost-oriented.  \Cref{tab:metrics}
summarizes key properties.

\begin{table}[t]
\centering
\footnotesize
\begin{tabular}{@{}llll@{}}
\toprule
\textbf{Metric} & \textbf{Category} & \textbf{Year} & \textbf{Models} \\
\midrule
McCabe CC    & Structural    & 1976 & Test paths \\
NPath        & Structural    & 1988 & Execution paths \\
ABC          & Structural    & 1997 & Assignment/Branch/Cond \\
Cognitive CC & Structural    & 2018 & Reading difficulty \\
Nesting Depth & Structural   & ---  & Stack depth \\
\midrule
SLOC         & Size          & ---  & Length \\
Token Count  & Size          & ---  & Lexical volume \\
\midrule
Halstead Vol & Vocabulary    & 1977 & Program vocabulary \\
Halstead Dif & Vocabulary    & 1977 & Implementation difficulty \\
\midrule
AIRD         & AI-cost       & 2025 & LLM reasoning effort \\
AICP         & AI-cost       & 2025 & LLM context tokens \\
\bottomrule
\end{tabular}
\caption{Complexity metrics surveyed, by category.}
\label{tab:metrics}
\end{table}

\subsection{Structural Metrics}

\subsubsection{McCabe Cyclomatic Complexity}

McCabe~\cite{mccabe} defines cyclomatic complexity as $V(G) = e - n + 2p$
for a control-flow graph with $e$ edges, $n$ nodes, and $p$ connected
components.  In practice this reduces to counting decision points plus
one: each \texttt{if}, \texttt{while}, \texttt{for}, \texttt{case},
\texttt{||}, and \texttt{\&\&} contributes one unit.

McCabe was designed to measure the minimum number of test cases needed to
achieve branch coverage.  It is the most widely deployed complexity metric:
tools, linters, and style guides from NASA's JPL~\cite{jpl} to MISRA
C~\cite{misra} cite McCabe thresholds (typically 10--15 as a warning
threshold, 25 as a hard limit).

\textbf{Limitation for AI cost:} McCabe counts branches but not nesting
depth.  A deeply nested function with 10 branches scores identically to a
flat function with 10 independent checks.  LLMs penalize nesting more
heavily than flatness: a nested structure forces the model to maintain a
context stack; a flat structure does not.

\subsubsection{NPath Complexity}

Nejmeh~\cite{npath} defines NPath as the number of acyclic execution paths
through a function.  For a function with $n$ independent branches,
NPath~$= 2^n$; two sequential \texttt{if} statements produce NPath~$= 4$.
NPath grows exponentially and is considered complementary to McCabe, which
grows linearly.

\textbf{Limitation:} NPath is rarely used in practice because exponential
growth makes thresholds hard to set and communicate.  The lizard tool
computes it but does not recommend thresholds.

\subsubsection{ABC Complexity}

Fitzpatrick~\cite{abc} defines ABC complexity as a three-component vector
$(A, B, C)$ where $A$ counts assignment operations, $B$ counts branch
operations (function calls and method dispatches), and $C$ counts condition
checks.  The scalar magnitude $\|ABC\| = \sqrt{A^2 + B^2 + C^2}$ is
reported as the complexity score.

ABC was designed for Ruby code quality and is the basis of Rails' Flog
tool.  It is useful when branch count and assignment density are both
relevant, but it conflates two orthogonal signals (function call density
and branching) into one scalar.

\subsubsection{Cognitive Complexity}

Campbell~\cite{cognitive-complexity} introduced Cognitive Complexity in
2018 as a direct response to the limitations of McCabe.  The key
differences are:

\begin{itemize}
    \item Nested structures incur a nesting penalty --- each level of
        nesting adds to the cost of the enclosing control structure.
    \item Sequences of \texttt{else if} are charged once, not once per
        branch.
    \item \texttt{switch} is a single increment regardless of arm count.
    \item Recursive calls and co-routines incur a flat increment.
\end{itemize}

Cognitive Complexity has been validated empirically: Campbell cites
SonarSource data showing higher correlation with developer-reported
difficulty than McCabe.  Mozilla's \texttt{rust-code-analysis} computes
Cognitive Complexity for Rust, C, and Python; knots was validated against
it at 1.004$\times$ mean ratio across 11,365 Rust functions.

\textbf{Limitation for AI cost:} Cognitive Complexity models \emph{human}
structural comprehension.  It does not account for external dependency
breadth (AICP) or state access patterns (state coupling).

\subsubsection{Nesting Depth}

Maximum nesting depth is a direct complement to Cognitive Complexity.
While Cognitive Complexity accumulates nesting penalties, nesting depth
reports the worst-case stack depth the reader must maintain.  A function
with a single 6-deep nesting chain scores low on McCabe and moderate on
Cognitive Complexity, but the nesting depth signal is unambiguous.

The JPL coding standard~\cite{jpl} limits nesting to 2 levels; MISRA C
recommends 4.  Both reflect the same intuition: deep nesting imposes
disproportionate maintenance cost.

\subsection{Size-Based Metrics}

\subsubsection{SLOC}

Source lines of code (SLOC) is the oldest and most broadly understood
complexity proxy.  It does not measure anything structural about the
code; it measures how long the reader must engage with it.  Long
functions are correlated with high complexity, but the correlation is
loose: a 300-line \texttt{switch} statement with 100 cases is long but
structurally simple.

Different tools define SLOC differently.  Physical SLOC counts every
non-blank line; logical SLOC counts statements.  Most modern tools use
physical SLOC minus blank lines and comment-only lines.  knots follows
this definition, verified against \texttt{cloc} counts.

\subsubsection{Token Count}

Lizard and some research tools report token count as a proxy for function
size.  Token count is more stable than SLOC across formatting styles
(brace placement, line wrapping) but is rarely used as a primary threshold
metric in practice.

\subsection{Vocabulary-Based Metrics}

\subsubsection{Halstead Software Science}

Halstead~\cite{halstead} proposed a family of metrics based on operator and
operand vocabulary:

\begin{align}
\text{Volume} &= N \cdot \log_2 \eta \\
\text{Difficulty} &= \frac{\eta_1}{2} \cdot \frac{N_2}{\eta_2} \\
\text{Effort} &= \text{Difficulty} \times \text{Volume}
\end{align}

where $\eta_1$ is the count of unique operators, $\eta_2$ unique operands,
$N_1$ total operator occurrences, $N_2$ total operand occurrences, and
$\eta = \eta_1 + \eta_2$.

Halstead metrics are computed by several tools (radon for Python,
rust-code-analysis for Rust) but have largely fallen out of favor in
applied software engineering.  Shepperd~\cite{shepperd} found that
Halstead's predictions were no better than SLOC alone for effort
estimation.  They remain in use in research as one component of composite
metrics (see maintainability index, below).

\subsection{Composite Metrics}

\subsubsection{Maintainability Index}

The Maintainability Index (MI)~\cite{mi} combines Halstead volume, McCabe
cyclomatic complexity, and SLOC into a single score:

\begin{align}
\text{MI} = 171 &- 5.2 \cdot \ln(\text{HV}) \nonumber \\
             &- 0.23 \cdot V(G) - 16.2 \cdot \ln(\text{SLOC})
\end{align}

MI is normalized to a 0--100 scale in some tools (Visual Studio uses a
variant).  It is the most widespread composite metric but suffers from the
compound weaknesses of its components: Halstead volume's weak predictive
value drags on the formula, and the logarithm compression makes the score
insensitive to large changes at the high end.

\subsection{AI-Cost Metrics}

AIRD and AICP are described in detail in the companion
paper~\cite{knots-aird}.  We include them in this taxonomy because they
represent a qualitatively new category: metrics calibrated to model
\emph{LLM cost} rather than human cognitive load.  The key design choices
--- separating reasoning effort (AIRD) from context pressure (AICP),
calibrating ceiling values against empirical p99 distributions, and
incorporating state coupling as a dampener for mechanical splits --- have
no precedent in the prior metrics literature.

%% ============================================================
\section{Tool Evaluation}
\label{sec:tools}

We evaluate 8 open-source tools against 5 dimensions: metric coverage,
language support, CI integration ergonomics, structured output formats,
and analysis speed.  \Cref{tab:tools} gives the overview; the subsections
provide detail on notable characteristics.

\begin{table*}[t]
\centering
\footnotesize
\setlength{\tabcolsep}{4pt}
\begin{tabular}{@{}lllllll@{}}
\toprule
\textbf{Tool} & \textbf{Lang(s)} & \textbf{Metrics} & \textbf{CI exit} & \textbf{JSON} & \textbf{pre-commit} & \textbf{Speed} \\
\midrule
pmccabe~\cite{pmccabe}              & C, C++           & CC                       & \checkmark & --         & --         & Fast \\
lizard~\cite{lizard}                & 27 langs         & CC, SLOC, params, tokens & \checkmark & \checkmark & --         & Fast \\
radon~\cite{radon}                  & Python only      & CC, Halstead, MI, SLOC   & \checkmark & \checkmark & --         & Fast \\
rust-code-analysis~\cite{rca}       & Rust, C, C++, JS & CC, Cognitive, Halstead, SLOC, many more & -- & \checkmark & -- & Fast \\
tokei~\cite{tokei}                  & 150+ langs       & SLOC only                & --         & \checkmark & --         & Very fast \\
SonarQube~\cite{sonarqube}          & 30+ langs        & CC, Cognitive, SLOC, duplications, smells & -- & REST API & -- & Slow (server) \\
understand~\cite{understand}        & 20+ langs        & CC, Halstead, many more  & --         & --         & --         & Moderate \\
knots~\cite{knots}                  & C, C++, Rust, Py, JS & CC, Cognitive, Nesting, SLOC, ABC, Returns, TestScore, ExternalCalls, AIRD, AICP, StateCoupling & \checkmark & \checkmark & \checkmark & Fast \\
\bottomrule
\end{tabular}
\caption{Open-source complexity tools surveyed.  ``CI exit'' = tool exits non-zero on threshold
violation.  ``pre-commit'' = native pre-commit framework integration.  ``Speed'' = qualitative
assessment on a 100K-LOC codebase.  CC = cyclomatic complexity.}
\label{tab:tools}
\end{table*}

\subsection{pmccabe}

pmccabe~\cite{pmccabe} is the canonical McCabe implementation for C and C++.
It has been in continuous use since the 1990s and produces output that
knots was validated against (100\% match across 32,205 functions).  Its
scope is deliberately narrow: it computes and reports McCabe complexity,
nothing else.  It has no structured output, no threshold configuration, and
no CI integration beyond shell scripting.

pmccabe remains the appropriate tool when the single concern is McCabe
compliance in a C/C++ shop, but it does not scale to the multi-metric,
multi-language, CI-native requirements of modern development.

\subsection{lizard}

Lizard~\cite{lizard} is a Python-based multi-language tool that computes
McCabe CC, SLOC, parameter count, and token count for 27 programming
languages.  It is the most broadly deployed open-source complexity scanner:
CI configurations referencing lizard appear in tens of thousands of public
repositories.

Lizard exits non-zero on threshold violations and produces JSON output,
making it the most CI-ready tool in the survey after knots.  Its primary
limitations are: it does not compute Cognitive Complexity; its SLOC
counting differs from standard definitions for several languages; and it
has no pre-commit framework integration, requiring manual CI configuration.

\subsection{radon}

Radon~\cite{radon} is a Python-only tool that computes the broadest metric
set of any single-language tool: McCabe CC, Halstead volume/difficulty/effort,
Maintainability Index, and SLOC.  It is the standard choice for Python
teams that want Halstead metrics without a heavyweight tool.  Its JSON
output and threshold-based exit codes make it CI-compatible.

The Python-only scope is radon's primary limitation.  Teams working in
polyglot codebases cannot use radon as a single-pane-of-glass solution.

\subsection{rust-code-analysis}

Mozilla's rust-code-analysis (RCA)~\cite{rca} computes the most comprehensive
metric set of any open-source tool: McCabe, Cognitive Complexity, Halstead
family, SLOC, comment density, function count, and several structural metrics.
It supports Rust, C/C++, Python, and JavaScript via tree-sitter.

RCA is a research artifact --- it was designed for Mozilla's internal code
quality research and as a validated reference implementation for Cognitive
Complexity.  It produces JSON output but does not provide threshold-based
exit codes, making CI integration purely manual.  There is no pre-commit
hook and no documented threshold recommendations.  The tool is not actively
maintained for deployment use cases.

\subsection{tokei}

Tokei~\cite{tokei} supports over 150 languages and reports SLOC, blank
lines, comment lines, and total lines per file and language.  It does not
compute any structural metric.  Its value is as a fast, reliable SLOC
counter across large polyglot codebases, not as a complexity analyzer.
We include it because many teams use tokei alongside a structural analyzer.

\subsection{SonarQube}

SonarQube~\cite{sonarqube} is a server-based platform that computes
McCabe CC, Cognitive Complexity, SLOC, code duplication, and code smells
across 30+ languages.  Its Cognitive Complexity implementation matches
Campbell's specification.  SonarQube is the most full-featured tool in
the survey, but its server deployment model, licensing structure (many
features require SonarQube Developer Edition or higher), and analysis
latency (full project scans take minutes to hours) make it unsuitable for
pre-commit or per-function CI gate use cases.

SonarCloud (the hosted SaaS variant) has better CI integration via GitHub
Actions / GitLab CI, but shares the same latency and cost characteristics.

\subsection{Understand}

Scientific Toolkit's Understand~\cite{understand} is a commercial static
analysis IDE that computes over 40 complexity metrics including McCabe,
Halstead, and several call-graph-based metrics.  It is the most
comprehensive tool in the survey in raw metric count, but it is
commercial, requires manual invocation through a GUI or script, and
does not integrate into pre-commit or lightweight CI workflows.

\subsection{knots}

Knots~\cite{knots} is a Rust-based, multi-language complexity analyzer
designed for CI-native deployment.  It computes 13 metrics per function:
McCabe CC, Cognitive Complexity, nesting depth, SLOC, ABC magnitude,
return count, test scoring (5 sub-dimensions), external call count,
AIRD, AICP, and state coupling.

Knots was validated against reference implementations: 100\% McCabe match
against pmccabe across 32,205 functions; 1.004$\times$ mean Cognitive
match against rust-code-analysis across 11,365 Rust functions.  It
supports C, C++, Rust, Python, and JavaScript via tree-sitter, exits
non-zero on configurable threshold violations, produces JSON/NDJSON/SARIF/CSV
output, and ships a pre-commit framework hook.

The metrics not found in any other surveyed tool are AIRD, AICP, and state
coupling, which are covered in the companion paper~\cite{knots-aird}.

%% ============================================================
\section{The AI-Cost Gap}
\label{sec:gap}

\subsection{What Existing Metrics Miss}

We identify three gaps between existing complexity metrics and the
requirements of AI-cost modeling.

\textbf{Context pressure is unmeasured.}  No surveyed metric captures the
number of external symbols a function requires resolved before it can be
understood.  A function that calls 15 external APIs forces an LLM to
expand 15 dependency chains before generating an edit; a 300-line linear
function forces tokenization of 300 lines.  These are distinct costs.
External call count captures the former; SLOC captures the latter.  No
existing composite metric combines them in a way calibrated to LLM token
budgets.

\textbf{State access patterns are unmeasured.}  Traditional metrics do not
distinguish between a function's explicit interface (parameter count) and
its actual state dependencies (fields accessed from a passed object or
\texttt{self}).  The state coupling term in AIRD was motivated by a
specific failure mode: a god-struct anti-pattern where struct-wrapping
parameters hides coupling from all arity-based metrics without actually
reducing the state the function touches
(\cref{sec:tools}).

\textbf{Reasoning vs.\ context are conflated.}  Existing composite metrics
(MI, Halstead Effort) combine structural and size-based signals into a
single score.  For AI cost, these are orthogonal: a function can be cheap
to load (low AICP) but hard to reason about (high AIRD), or vice versa.
Separating them gives the developer two actionable levers: reduce external
calls to lower AICP, reduce cognitive complexity to lower AIRD.

\subsection{What Existing Tools Miss}

Beyond metrics, existing tools have deployment gaps relevant to AI-cost
use cases.

\textbf{Pre-commit integration.}  Of the 8 surveyed tools, only knots
ships a pre-commit framework hook.  The others require manual CI
configuration or shell scripting.  Pre-commit is the appropriate gate
for per-function complexity enforcement because it fires at the smallest
unit of change (a commit), not at the coarse granularity of a CI pipeline
run.

\textbf{Multi-language support with per-function granularity.}  Most
tools are either single-language (radon, pmccabe) or multi-language but
file-granularity (tokei).  Lizard offers multi-language per-function
output but lacks Cognitive Complexity and AI-cost metrics.

\textbf{Structured output for downstream tooling.}  SARIF
output~\cite{sarif} is increasingly required for integration with GitHub
Code Scanning and similar platforms.  Of the surveyed tools, only SonarQube
and knots produce SARIF; lizard and radon produce JSON but not SARIF.

\subsection{Threshold Calibration}

Threshold recommendations in the literature vary widely and are rarely
empirically grounded.  McCabe thresholds range from 10 (JPL~\cite{jpl})
to 25 (many style guides) to ``no limit'' (some Agile practitioners who
argue test coverage is a better gate).  No tool other than knots documents
the corpus and distribution from which its default thresholds were derived.

The AIRD threshold of 85 was derived from the p99 distribution across
32,205 functions from 6 production codebases, then empirically validated
against Sonnet 4.6 and Opus 4.8 behavior.  This calibration methodology
--- corpus-derived ceilings with LLM-behavior validation --- is
a contribution independent of the specific metric, and one not found in
prior tool documentation.

%% ============================================================
\section{Related Work}
\label{sec:related}

\textbf{Complexity metric surveys.}
Zuse~\cite{zuse} provides a comprehensive theoretical treatment of software
complexity measures, covering axiomatic properties and measurement-theoretic
validity.  Shepperd~\cite{shepperd} conducts an early empirical comparison
of McCabe and Halstead against simpler size metrics, finding that SLOC
predicts defect density as well as more sophisticated metrics.  Gill and
Kemerer~\cite{gill-kemerer} compare 8 metrics on a commercial COBOL
codebase; they find high inter-correlation, suggesting most metrics capture
the same underlying construct.

\textbf{Cognitive Complexity.}
Campbell~\cite{cognitive-complexity} is the original specification.
Gruner~\cite{gruner} provides a formal treatment of why Cognitive
Complexity better predicts developer comprehension time than McCabe.

\textbf{AI code generation quality.}
Studies of LLM code generation quality (e.g.,
Chen~et~al.~\cite{chen-codex} on Codex) evaluate correctness on benchmark
problems but do not study how input function complexity affects LLM
modification accuracy.  This is the empirical gap motivating the
companion paper's AIRD/AICP validation work.

\textbf{Technical debt metrics.}
SQALE~\cite{sqale} and the Maintainability Index~\cite{mi} aggregate
multiple metrics into a single remediation cost estimate.  Neither was
designed for LLM cost estimation, but the composite-metric approach
informed AIRD's design.

%% ============================================================
\section{Conclusion}
\label{sec:conclusion}

We have surveyed 10 complexity metrics and 8 open-source tools, characterizing
where existing work falls short for AI-cost-aware complexity analysis.  The
gaps are consistent: no existing metric captures LLM context pressure or
state access patterns; no tool other than knots ships both multi-language
per-function analysis and pre-commit CI integration; and no tool documents
empirically calibrated thresholds derived from a production codebase corpus.

The AIRD and AICP metrics described in the companion paper address the
first gap directly.  We release knots as an open-source tool to address
the deployment gap, with a pre-commit hook, SARIF output, and documented
threshold recommendations derived from a 32,205-function corpus.

\bibliographystyle{plainnat}

\begin{thebibliography}{99}

\bibitem{mccabe}
T.~J. McCabe.
\newblock A complexity measure.
\newblock \emph{IEEE Transactions on Software Engineering}, 2(4):308--320,
  1976.

\bibitem{halstead}
M.~H. Halstead.
\newblock \emph{Elements of Software Science}.
\newblock Elsevier North-Holland, New York, 1977.

\bibitem{cognitive-complexity}
G.~A. Campbell.
\newblock {Cognitive Complexity: A new way of measuring understandability}.
\newblock SonarSource SA, Geneva, Switzerland, 2018.
\newblock \url{https://www.sonarsource.com/resources/cognitive-complexity/}.

\bibitem{npath}
B.~A. Nejmeh.
\newblock {NPath}: A measure of execution path complexity and its applications.
\newblock \emph{Communications of the ACM}, 31(2):188--200, 1988.

\bibitem{abc}
J.~Fitzpatrick.
\newblock Applying the {ABC} metric to {C}, {C++} and {Java}.
\newblock \emph{C++ Report}, 9(9), 1997.

\bibitem{mi}
P.~Oman and J.~Hagemeister.
\newblock Metrics for assessing a software system's maintainability.
\newblock In \emph{Proc.\ IEEE Conf.\ on Software Maintenance}, pages
  337--344, 1992.

\bibitem{pmccabe}
B.~Raiter.
\newblock {pmccabe} --- McCabe cyclomatic complexity tool for C and C++.
\newblock \url{https://people.debian.org/~bame/pmccabe/}, 1994.

\bibitem{lizard}
T.~Yin.
\newblock {Lizard}: A code complexity analyzer.
\newblock \url{https://github.com/terryyin/lizard}, 2013.

\bibitem{radon}
M.~Barzoni.
\newblock {Radon}: Code metrics for Python.
\newblock \url{https://radon.readthedocs.io/}, 2013.

\bibitem{rca}
{Mozilla}.
\newblock {rust-code-analysis}: A library to analyze and extract information
  from source code.
\newblock \url{https://mozilla.github.io/rust-code-analysis/}, 2020.

\bibitem{tokei}
X.~Kipp.
\newblock {Tokei}: A program that displays statistics about your code.
\newblock \url{https://github.com/XAMPPRocky/tokei}, 2016.

\bibitem{sonarqube}
{SonarSource}.
\newblock {SonarQube}: Code quality and security.
\newblock \url{https://www.sonarqube.org/}, 2007.

\bibitem{understand}
{Scientific Toolkit}.
\newblock {Understand}: Static analysis and metrics.
\newblock \url{https://www.scitools.com/}, 1993.

\bibitem{knots}
B.~Arrendondo.
\newblock {knots}: A multi-language code complexity analyzer with AI cost
  metrics.
\newblock \url{https://github.com/brandon-arrendondo/knots}, 2024.

\bibitem{knots-aird}
B.~Arrendondo.
\newblock {AIRD and AICP: Metrics for predicting AI reasoning and context cost
  in code modification}.
\newblock % TODO: journal/conference, year

\bibitem{jpl}
G.~J. Holzmann.
\newblock {The Power of 10: Rules for Developing Safety-Critical Code}.
\newblock \emph{IEEE Computer}, 39(6):95--97, 2006.

\bibitem{misra}
{MISRA Consortium}.
\newblock {MISRA C: 2012 --- Guidelines for the Use of the C Language in
  Critical Systems}, 2012.

\bibitem{sarif}
{OASIS SARIF Technical Committee}.
\newblock {Static Analysis Results Interchange Format (SARIF) Version 2.1.0}.
\newblock OASIS Standard, 2020.
\newblock \url{https://docs.oasis-open.org/sarif/sarif/v2.1.0/sarif-v2.1.0.html}.

\bibitem{zuse}
H.~Zuse.
\newblock \emph{A Framework of Software Measurement}.
\newblock Walter de Gruyter, Berlin, 1998.

\bibitem{shepperd}
M.~Shepperd.
\newblock A critique of cyclomatic complexity as a software metric.
\newblock \emph{Software Engineering Journal}, 3(2):30--36, 1988.

\bibitem{gill-kemerer}
G.~K. Gill and C.~F. Kemerer.
\newblock Cyclomatic complexity density and software maintenance productivity.
\newblock \emph{IEEE Transactions on Software Engineering}, 17(12):1284--1288,
  1991.

\bibitem{gruner}
S.~Gruner.
\newblock On the conceptual difference between cognitive and cyclomatic
  complexity.
\newblock \emph{South African Computer Journal}, 54:53--71, 2014.

\bibitem{chen-codex}
M.~Chen, J.~Tworek, H.~Jun, Q.~Yuan, H.~P. de~Oliveira~Pinto, J.~Kaplan,
  H.~Edwards, Y.~Burda, N.~Joseph, G.~Brockman, et~al.
\newblock Evaluating large language models trained on code.
\newblock \emph{arXiv preprint arXiv:2107.03374}, 2021.

\bibitem{sqale}
J.-L. Letouzey.
\newblock {The SQALE Method for Evaluating Technical Debt}.
\newblock In \emph{Proc.\ 3rd Intl.\ Workshop on Managing Technical Debt},
  pages 31--36, 2012.

\end{thebibliography}

\end{document}
